\chapter{Introduzione: AI simbolica e bisogno della logica}

\textit{Modulo 2 -- M. Gabbrielli, DISI, Università di Bologna.}
Lucidi: \texttt{lucidi/1.Intro.Logica}. Materiale di approfondimento:
Van Dalen, \emph{Logic and Structure}; Bratko, \emph{Prolog Programming for AI};
Gabbrielli--Martini, capitolo su LP e CLP; Tutorial MiniZinc;
Russell \& Norvig, cap.~1.

% =====================================================================
\section{Organizzazione del corso}
% =====================================================================

\nt{
  \begin{itemize}
    \item 1\textsuperscript{o} modulo: Felli (sulla slide il prof ha corretto
          a mano ``Lodi'' $\to$ ``Felli'').
    \item 2\textsuperscript{o} modulo: questo (Gabbrielli).
    \item Il corso integrato con Dal Lago è stato \textbf{cancellato} a mano dal
          prof: non fa parte di questo esame.
    \item Esami \textbf{separati} per ogni modulo; per questo modulo l'esame è
          \textbf{solo scritto}.
    \item Parte del materiale viene dal corso di logica (in italiano) di
          \textbf{Claudio Sacerdoti Coen} (UniBo), nome scritto a mano sulla
          copertina; il resto da Russell--Norvig, Van Dalen, E.~Clarke.
  \end{itemize}
}

\subsection{Programma del modulo}
\begin{enumerate}
  \item AI simbolica e bisogno della logica.
  \item Introduzione alla logica: logica proposizionale, logica del primo
        ordine (FOL), risoluzione, unificazione.
  \item Logic Programming: interpretazione procedurale, interpretazione
        dichiarativa, Prolog (prima parte; la seconda in un altro corso).
  \item Constraint Programming: CLP, MiniZinc (prima parte).
\end{enumerate}

\subsection{Testi}
\begin{itemize}
  \item Dispense (su Virtuale).
  \item D.~van Dalen, \emph{Logic and Structure}, 4\textsuperscript{a} ed.,
        Springer --- riferimento per la parte di logica.
  \item I.~Bratko, \emph{Prolog Programming for Artificial Intelligence}.
  \item M.~Triska, \emph{The Power of Prolog} (online).
  \item Russell, Norvig, \emph{Artificial Intelligence: A Modern Approach}.
  \item \emph{The MiniZinc Handbook} (online).
\end{itemize}

% =====================================================================
\section{Due esempi motivazionali}
% =====================================================================

Prima di parlare di logica il prof mostra due problemi che si risolvono
\emph{descrivendoli} invece di scrivere un algoritmo:
\begin{enumerate}
  \item un problema \textbf{combinatorio} (risolto in Prolog);
  \item un problema di \textbf{ottimizzazione} (risolto in MiniZinc).
\end{enumerate}

\subsection{Esempio 1: un problema combinatorio}

\dfn{Il problema (sequenza di Langford a tre copie)}{
  Disporre in una sequenza tre $1$, tre $2$, \dots, tre $9$ (in totale
  $27$ numeri) in modo che, per ogni $i \in [1,9]$, tra due occorrenze
  \emph{successive} di $i$ ci siano \textbf{esattamente $i$ numeri}.
}

\ex{Una soluzione}{
  \[
    1,9,1,2,1,8,2,4,6,2,7,9,4,5,8,6,3,4,7,5,3,9,6,8,3,5,7
  \]
  Verifica per $i=1$: gli $1$ stanno nelle posizioni $1,3,5$; tra due $1$
  consecutivi c'è esattamente un numero. Per $i=2$: posizioni $4,7,10$,
  quindi due numeri in mezzo. (Sulla slide il prof sottolinea a colori le
  occorrenze di ciascun numero.)
}

Il programma Prolog che lo risolve:

\begin{lstlisting}[language=Prolog]
sequence([_,_,_,_,_,_,_,_,_,_,_,_,_,_,_,_,_,_,_,_,_,_,_,_,_,_,_]).
% sequence(X) -> X e' una lista di 27 variabili

question(S) :-
    sequence(S),
    sublist([9,_,_,_,_,_,_,_,_,_,9,_,_,_,_,_,_,_,_,_,9], S),
    sublist([8,_,_,_,_,_,_,_,_,8,_,_,_,_,_,_,_,_,8], S),
    sublist([7,_,_,_,_,_,_,_,7,_,_,_,_,_,_,_,7], S),
    sublist([6,_,_,_,_,_,_,6,_,_,_,_,_,_,6], S),
    sublist([5,_,_,_,_,_,5,_,_,_,_,_,5], S),
    sublist([4,_,_,_,_,4,_,_,_,_,4], S),
    sublist([3,_,_,_,3,_,_,_,3], S),
    sublist([2,_,_,2,_,_,2], S),
    sublist([1,_,1,_,1], S).
% S e' una soluzione
\end{lstlisting}

\nt{
  \textbf{Come leggerlo (approfondimento, Bratko cap.~1--3).}
  \begin{itemize}
    \item \texttt{\_} è una \textbf{variabile anonima}: ogni occorrenza è una
          variabile nuova, diversa dalle altre.
    \item Il programma \emph{non dice come} cercare la soluzione: dice solo
          \emph{che cosa} è una soluzione. $S$ è una lista di 27 elementi, e per
          ogni $i$ il ``pattern'' $[i, \underbrace{\_,\dots,\_}_{i}, i,
          \underbrace{\_,\dots,\_}_{i}, i]$ deve comparire come
          \textbf{sottolista contigua} di $S$.
    \item \texttt{sublist/2} non è un predicato standard di SWI-Prolog; una
          definizione possibile è
          \begin{center}
          \texttt{sublist(L, S) :- append(\_, T, S), append(L, \_, T).}
          \end{center}
          cioè $L$ è sottolista di $S$ se $S = A \cdot L \cdot B$ per qualche
          $A,B$.
    \item Il ``come'' lo fa l'interprete: \textbf{unificazione} (il pattern
          viene ``incollato'' su un tratto di $S$, legando le variabili) e
          \textbf{backtracking} (se un \texttt{sublist} successivo fallisce si
          torna indietro e si prova un'altra posizione).
          Mettere prima il $9$ (il pattern più lungo) riduce le posizioni
          possibili e quindi lo spazio di ricerca.
  \end{itemize}
}

\subsection{Esempio 2: un problema di ottimizzazione (Cakes)}

\dfn{Il problema delle torte}{
  Si possono fare due tipi di torte:
  \begin{enumerate}
    \item \textbf{banana}: 250\,g farina, 2 banane, 75\,g zucchero, 100\,g burro;
    \item \textbf{cioccolato}: 200\,g farina, 75\,g cacao, 150\,g zucchero,
          150\,g burro.
  \end{enumerate}
  Prezzi di vendita: cioccolato \$4.50, banana \$4.00.
  Disponibilità: 4\,kg farina, 6 banane, 2\,kg zucchero, 500\,g burro,
  500\,g cacao.
  \textbf{Domanda}: quante torte di ciascun tipo fare per \textbf{massimizzare
  il profitto}?
}

\begin{lstlisting}[language=MiniZinc]
% Baking cakes for the school fete
var 0..100: b;   % no. of banana cakes
var 0..100: c;   % no. of chocolate cakes

constraint 250*b + 200*c <= 4000;   % flour
constraint 2*b <= 6;                % bananas
constraint 75*b + 150*c <= 2000;    % sugar
constraint 100*b + 150*c <= 500;    % butter
constraint 75*c <= 500;             % cocoa

solve maximize 400*b + 450*c;       % profit (in cents)

output ["no. of banana cakes = \(b)\n",
        "no. of chocolate cakes = \(c)\n"];
\end{lstlisting}

\nt{
  \textbf{Struttura di un modello MiniZinc (Tutorial MiniZinc).}
  \begin{itemize}
    \item \texttt{var 0..100: b;} dichiara una \textbf{variabile decisionale}
          intera con dominio $[0,100]$;
    \item \texttt{constraint} aggiunge un \textbf{vincolo} che ogni soluzione
          deve rispettare;
    \item \texttt{solve maximize} $f$ chiede una soluzione \textbf{ottima}
          rispetto a $f$ (con \texttt{solve satisfy} se ne chiede una qualsiasi);
    \item i prezzi sono scritti in centesimi per restare negli interi.
  \end{itemize}
  Anche qui scriviamo solo \emph{che cosa} vogliamo; il solver decide come
  esplorare le combinazioni.
}

\ex{La soluzione ottima (approfondimento)}{
  Il vincolo stringente è il burro: $100b + 150c \le 500$, cioè
  $2b + 3c \le 10$; inoltre le banane danno $b \le 3$. Le coppie
  ammissibili ``sul bordo'' sono:
  \[
  \begin{array}{c|c|c}
    (b,c) & \text{burro} & 400b + 450c \\ \hline
    (3,1) & 450 & 1650 \\
    (2,2) & 500 & \mathbf{1700} \\
    (1,2) & 400 & 1300 \\
    (0,3) & 450 & 1350
  \end{array}
  \]
  Gli altri vincoli (farina, zucchero, cacao) sono soddisfatti da tutte.
  L'ottimo è $b = 2$, $c = 2$, profitto \$17.00.
}

% =====================================================================
\section{Programmazione dichiarativa}
% =====================================================================

\dfn{Algorithm = Logic + Control (Kowalski)}{
  Un algoritmo si scompone in due parti:
  \begin{itemize}
    \item \textbf{Logic}: \emph{che cosa} va calcolato (la specifica del problema);
    \item \textbf{Control}: \emph{come} lo si calcola (ordine dei passi,
          strategia di ricerca).
  \end{itemize}
}

\dfn{Programmazione tradizionale vs.\ dichiarativa}{
  \begin{itemize}
    \item Nella programmazione \textbf{tradizionale} (imperativa) il
          programmatore si occupa sia della Logic sia del Control.
    \item Nella programmazione \textbf{dichiarativa} il programmatore si
          occupa \textbf{solo della Logic}; il Control lo gestisce
          l'interprete del linguaggio.
  \end{itemize}
  Il compito del programmatore è \emph{dichiarare} il problema: \textbf{cosa}
  vogliamo ottenere e non \textbf{come} ottenerlo.
}

\dfn{Tre paradigmi dichiarativi}{
  \begin{enumerate}
    \item \textbf{Logic Programming} (Prolog, \dots): basato sulla nozione di
          \emph{relazione}.
    \item \textbf{Functional programming} (ML, Haskell, OCaml, \dots): basato
          sulla composizione di \emph{funzioni}.
    \item \textbf{Constraint Programming} (MiniZinc, CLP, ILOG, \dots): basato
          su relazioni e su \emph{solver} specifici per un dominio.
  \end{enumerate}
}

\nt{
  \emph{``Constraint programming represents one of the closest approaches
  computer science has yet made to the Holy Grail of programming: the user
  states the problem, the computer solves it.''} --- E.~C.~Freuder.

  \textbf{Approfondimento (Gabbrielli--Martini).} In Prolog il Control è
  fissato: le clausole si provano dall'alto in basso, i letterali del corpo
  da sinistra a destra, con ricerca in profondità e backtracking.
  Per questo in pratica l'ordine delle clausole \emph{conta} (può cambiare
  efficienza e perfino terminazione), e la separazione Logic/Control in
  Prolog non è perfetta. Il significato dichiarativo (l'insieme delle
  risposte \emph{vere}) resta però indipendente da quest'ordine.
}

% =====================================================================
\section{Che cos'è la logica}
% =====================================================================

\dfn{Logica}{
  La logica è lo studio sistematico della \textbf{forma delle inferenze
  valide} e delle leggi più generali della \textbf{verità}.
  Un'\textbf{inferenza valida} è una in cui esiste una specifica relazione di
  supporto logico tra le \emph{premesse} e la \emph{conclusione}.
}

Tradizionalmente la logica comprende: classificazione degli argomenti,
esposizione della ``forma logica'' comune a tutti gli argomenti validi, studio
della dimostrazione e dell'inferenza (compresi paradossi e fallacie), studio di
\textbf{sintassi} e \textbf{semantica}. È studiata in filosofia (dall'antichità),
in matematica (da metà Ottocento) e oggi nelle scienze cognitive (informatica,
linguistica, filosofia, psicologia).

\subsection{Tante logiche diverse}

Ogni logica fissa \emph{che cosa} conserva un'inferenza. Con premesse $X, Y$ e
conclusione $W$:

\begin{center}
\begin{tabular}{lll}
\toprule
\textbf{Logica} & \textbf{Domanda} & \textbf{Uso in CS/AI} \\
\midrule
Classica (prop., FOL) & se $X$ e $Y$ sono veri, $W$ è vero? & AI, linguaggi logici, verifica \\
Intuizionista & da prove di $X$, $Y$ costruisco una prova di $W$? & software corretto per costruzione \\
Lineare & avendo $X$ e $Y$ posso \emph{scambiarli} per $W$? & controllo delle risorse \\
Epistemica & se credo che tu sappia $X$, $Y$, credo $W$? & sicurezza \\
Temporale & se $X$, $Y$ accadranno, $W$ potrà accadere? & specifica e verifica \\
\bottomrule
\end{tabular}
\end{center}

\nt{
  Logiche diverse, ma problemi e tecniche simili. La logica è
  interdisciplinare (filosofia, matematica, informatica, AI).
  \textbf{Logica matematica} = applicazione del metodo matematico allo studio
  della logica, \emph{ma anche} logica come fondamento della matematica
  (logiche diverse $\Rightarrow$ matematiche diverse).
}

% =====================================================================
\section{Perché studiare logica in informatica}
% =====================================================================

Cinque motivi (numerati a mano dal prof sulla slide):
\begin{enumerate}
  \item l'informatica deriva dalla logica;
  \item la logica fonda la computazione e i linguaggi dichiarativi;
  \item la logica fonda l'informatica teorica (calcolabilità);
  \item serve a dimostrare la correttezza del software;
  \item la logica è un linguaggio privo di paradossi.
\end{enumerate}

\subsection{1) L'informatica deriva dalla logica}

\qs{È nato prima il computer o l'informatica?}{
  La risposta sembra ovvia, ma l'informatica esiste \textbf{prima e
  indipendentemente} dal computer, come \textbf{studio della calcolabilità}.
}

La storia, in breve:
\begin{itemize}
  \item La matematica antica (geometria a parte) non aveva basi formali
        rigorose; queste vengono cercate nella logica e nell'aritmetica.
  \item Si applicano metodi matematici allo studio della logica: la
        derivazione logica viene identificata con un processo di calcolo
        (\emph{calculemus!}).
  \item Studiando i processi di calcolo si scoprono risultati sorprendenti:
        esistono funzioni \textbf{non calcolabili} (anzi, sono la maggioranza),
        ed esiste una classe di funzioni calcolabili definibile in modi del
        tutto diversi (lambda calcolo, macchine di Turing, programmi logici,
        RAM).
  \item Intorno al 1930 nasce l'informatica.
\end{itemize}

\nt{
  \emph{``Computer Science is no more about computers than Astronomy is about
  telescopes.''} --- E.~W.~Dijkstra.

  Il fatto che formalismi così diversi definiscano \emph{la stessa} classe di
  funzioni è alla base della \textbf{tesi di Church--Turing}: le funzioni
  ``calcolabili da un procedimento effettivo'' sono esattamente quelle
  calcolabili da una macchina di Turing.
}

\subsection{2) Fondamento della computazione e dei linguaggi dichiarativi}

Due grandi classi di linguaggi: \textbf{dichiarativi} (funzionali, logici, a
vincoli) e \textbf{imperativi}.
\begin{itemize}
  \item \textbf{Linguaggi funzionali} (LISP, ML, Haskell, OCaml): basati sulla
        composizione di funzioni. Legame forte con la logica:
        \emph{programmare corrisponde a dimostrare} (in una logica
        opportuna) e \emph{i tipi corrispondono ad enunciati di teoremi}
        (corrispondenza di Curry--Howard).
  \item \textbf{Linguaggi logici} (Prolog, CLP, CHR, CCP): basati sulla nozione
        di relazione. Legame ancora più forte: \textbf{computazione =
        deduzione}; Algorithm = Logic + Control.
  \item \textbf{Linguaggi a vincoli} (CLP, CHR, CCP, MiniZinc): relazioni +
        solver specifici di dominio; esistono sia basati sulla logica sia
        imperativi; molto usati in applicazioni pratiche.
\end{itemize}

\subsection{3) Fondamento dell'informatica teorica (calcolabilità)}

Due domande:
\begin{enumerate}
  \item Data un'affermazione vera o falsa, possiamo sempre \emph{dimostrarla}?
  \item Data una funzione generica, possiamo \emph{calcolarla} con un computer?
\end{enumerate}
La risposta è \textbf{no} in entrambi i casi: la logica studia i limiti della
conoscenza. Non tutto si può dimostrare vero o falso (incompletezza), non
tutto si può calcolare. Sono limiti \emph{inevitabili} per l'informatica e
l'AI.

\subsection{4) Dimostrare la correttezza del software}

\qs{Possiamo essere sicuri che un generico programma scritto in un linguaggio
comune sia corretto?}{}

\thm{Risultato fondamentale}{
  La maggior parte delle proprietà (non banali) dei programmi \textbf{non è
  decidibile}. In particolare non può esistere uno strumento automatico che,
  dato un programma generico (in un linguaggio Turing-completo), decida se è
  corretto (ad esempio se termina).
}

\nt{
  Questo è il \textbf{teorema di Rice} (approfondimento): ogni proprietà non
  banale della \emph{funzione calcolata} da un programma è indecidibile.
  Il caso della terminazione è dimostrato più sotto
  (Teorema~\ref{th:halting}).
}

Come se ne esce:
\begin{itemize}
  \item si scrive una \textbf{specifica} e si usa la logica per dimostrare
        che il programma la soddisfa;
  \item \textbf{prove automatiche}: solo in casi specifici o su
        \emph{approssimazioni} delle proprietà;
  \item \textbf{prove assistite}: l'umano fa la dimostrazione (in logica), il
        computer la \emph{verifica}.
\end{itemize}

\nt{
  \textbf{Note a mano del prof: il disegno della ``trusted code base''.}
  \begin{itemize}
    \item Dato un programma (il foglio) e un oracolo ``?'' che dovrebbe dire
          se è corretto: un tale oracolo è \textbf{non decidibile} (croce
          rossa).
    \item Strada alternativa: l'umano scrive la \emph{specifica} e la
          \emph{dimostrazione} in logica (il riquadro con $\forall, \exists,
          \Rightarrow$); un piccolo programma \textbf{verificatore di prove}
          controlla la dimostrazione e risponde sì/no (spunta verde).
    \item Controllare una prova data è \emph{decidibile}, trovarla in generale
          no. Ci si deve fidare solo del verificatore, che è piccolo e
          semplice: questa è la \textbf{trusted code base} (evidenziata in
          giallo). È l'idea dei proof assistant (Coq/Rocq, Lean, Isabelle).
  \end{itemize}
}

\dfn{Specifica di un algoritmo di ordinamento (dagli appunti del prof)}{
  \textbf{Pre-condizione:} l'input è un array $b$ di interi.

  \textbf{Post-condizione:} l'output è un array $a$ di interi tale che
  \[
    \bigl(\forall i < a.\mathit{length}-1.\ a[i] \le a[i+1]\bigr)
    \ \wedge\ \bigl(\forall i < a.\mathit{length}.\ \exists j.\ a[i] = b[j]\bigr)
    \ \wedge\ \bigl(\forall i < b.\mathit{length}.\ \exists j.\ b[i] = a[j]\bigr).
  \]
  La prima congiunzione dice che $a$ è \textbf{ordinato}; le altre due che $a$
  e $b$ contengono \textbf{gli stessi valori}.
}

\nt{
  \textbf{Osservazione (approfondimento).} Le ultime due clausole dicono
  ``stessi valori'' ma \emph{non} ``stesse molteplicità'': l'output
  $a = [1,2,2]$ soddisfa la specifica con input $b = [2,1,1]$.
  La specifica corretta richiede che $a$ sia una \textbf{permutazione} di $b$
  (stessa lunghezza e, per ogni valore $x$, lo stesso numero di occorrenze).
  Esempio di quanto è facile sbagliare una specifica: anche la specifica va
  verificata. (Negli appunti gli indici sono scritti con $\le$; con array
  indicizzati da $0$ servono le disuguaglianze strette, come sopra.)
}

\subsubsection{Tre limiti delle funzioni calcolate da un computer}

\textbf{Nota importante:} i risultati seguenti valgono per linguaggi
\textbf{Turing-completi}, cioè con il massimo potere espressivo noto.

\mprop{Limite 1: le funzioni calcolate non sono solo totali}{
  Definiamo nel linguaggio
  \[ f(g) = \mathtt{not}\,(g(g)) \]
  e applichiamo $f$ a sé stessa: $f(f) = \mathtt{not}\,(f(f))$.
  Se $f(f)$ restituisse un booleano $v$ avremmo $v = \neg v$, assurdo.
  Quindi la computazione di $f(f)$ \textbf{non termina}: le funzioni calcolate
  dai programmi sono in generale \textbf{parziali}.
}

\nt{
  Accanto a questa slide il prof ha scritto
  \[ X = \{Y \mid Y \notin Y\} \qquad X \in X \iff X \notin X \]
  cioè il \textbf{paradosso di Russell} (Sezione~\ref{sec:paradossi}):
  $f(g) = \mathtt{not}(g(g))$ è la stessa costruzione, con
  ``$g$ applicata a $g$'' al posto di ``$Y \in Y$''. In matematica genera un
  paradosso; in un linguaggio di programmazione genera una funzione
  \emph{parziale}.
}

\begin{Theorem}{Limite 2: la terminazione non è decidibile}{halting}
  Non esiste un programma $f$ (in alcun linguaggio) tale che, per ogni
  programma $g$ e argomento $x$,
  \[
    f(g,x) = \begin{cases}
      \mathtt{true} & \text{se } g(x) \text{ termina},\\
      \mathtt{false} & \text{se } g(x) \text{ non termina}.
    \end{cases}
  \]
\end{Theorem}

\pf{Dimostrazione}{
  Per assurdo, supponiamo che $f$ esista. Definiamo
  \begin{center}
  \texttt{h(g) = if f(g,g) = true then while true do skip else return 0}
  \end{center}
  e consideriamo $h(h)$:
  \[
    h(h) \text{ termina} \iff f(h,h) = \mathtt{false}
    \iff h(h) \text{ non termina}.
  \]
  Contraddizione, quindi $f$ non può esistere.
}

\nt{
  $f$ deve essere essa stessa \emph{totale} (risponde sempre) perché
  l'argomento funzioni: è proprio questo che la rende impossibile.
  La struttura è di nuovo la \textbf{diagonalizzazione}: $h$ applicata a sé
  stessa fa il contrario di quello che $f$ prevede.
}

\mlenma{I programmi sono numerabili}{
  L'insieme dei programmi esprimibili in un qualsiasi linguaggio di
  programmazione è \textbf{numerabile}.
}

\pf{Dimostrazione}{
  Un programma è una stringa finita su un alfabeto finito $\Sigma$. Le
  stringhe di lunghezza $n$ sono $|\Sigma|^n$, un numero finito. Elencando
  prima le stringhe di lunghezza $0$, poi quelle di lunghezza $1$, poi $2$,
  \dots\ (ciascun gruppo in ordine lessicografico) si ottiene
  un'enumerazione di $\Sigma^*$. I programmi sono un sottoinsieme di
  $\Sigma^*$, quindi sono al più numerabili.
}

\thm{Le funzioni booleane non sono numerabili}{
  Sia $F = \{ f \mid f : \mathbb{N} \to \{0,1\} \}$. Allora $F$ non è
  numerabile.
}

\pf{Dimostrazione (diagonale)}{
  Per assurdo, sia $f_1, f_2, f_3, \dots$ un'enumerazione di tutte le
  funzioni di $F$. Definiamo
  \[
    f^*(n) = \begin{cases} 1 & \text{se } f_n(n) = 0,\\ 0 & \text{se } f_n(n) = 1. \end{cases}
  \]
  Allora $f^* \in F$, ma per ogni $i$ si ha $f^*(i) \ne f_i(i)$, quindi
  $f^* \ne f_i$ per ogni $i$. Dunque $f^*$ non compare nell'enumerazione,
  contro l'ipotesi che fosse completa.
}

\cor{Limite 3: i linguaggi esprimono pochissime funzioni}{
  Nessun linguaggio di programmazione può esprimere tutte le funzioni
  booleane $\mathbb{N}\to\{0,1\}$ (e quindi, a fortiori, tutte le funzioni
  matematiche).
}

\pf{Dimostrazione}{
  I programmi sono numerabili (lemma), $F$ no (teorema): non può esistere una
  suriezione dai programmi su $F$.
}

\subsubsection{Dimostrazione alternativa: il diagonale di Cantor}

\thm{Cantor}{
  Sia $T$ un insieme con almeno due elementi. Allora $|T| < |T^T|$, dove
  $T^T$ è l'insieme delle funzioni $T \to T$.
}

\pf{Dimostrazione}{
  Non può essere $|T| > |T^T|$: la mappa $t \mapsto (\text{funzione costante }
  t)$ è iniettiva da $T$ in $T^T$. Basta dunque mostrare $|T| \ne |T^T|$.

  Per assurdo, sia $g : T \to T^T$ una biiezione. Definiamo $f \in T^T$
  ponendo, per ogni $t \in T$, $f(t) = t'$ con $t' \ne g(t)(t)$
  (esiste perché $T$ ha almeno due elementi). Allora per ogni $t$ si ha
  $f(t) \ne g(t)(t)$, quindi $f \ne g(t)$. Dunque $f$ non è nell'immagine
  di $g$, che non è suriettiva: assurdo.
}

\cor{Fatto 3, versione alternativa}{
  Sia $T$ l'insieme di tutti i valori di un linguaggio, che contiene almeno
  i booleani e le funzioni. Se le funzioni del linguaggio fossero
  \emph{tutte} le funzioni matematiche $T\to T$, avremmo
  $T \supseteq \{0,1\} \cup T^T$, quindi $|T| \ge |T^T|$, contro Cantor.
  Quindi un linguaggio non può esprimere tutte le funzioni matematiche.
}

\subsection{5) La logica come linguaggio privo di paradossi}
\label{sec:paradossi}

\dfn{Paradosso}{
  Un \textbf{paradosso} è un ragionamento apparentemente corretto (rispetto
  a regole e assiomi dati), da premesse apparentemente corrette, che porta a
  una conclusione apparentemente contraddittoria.
}

Il linguaggio naturale permette di esprimere paradossi, quindi non è adatto
come base per la matematica o l'informatica. La logica fornisce linguaggi
privi di paradossi.

\ex{Falso paradosso (trova l'errore)}{
  \[
    x = 1 \Rightarrow x^2 = x \Rightarrow x^2 - 1 = x - 1
    \Rightarrow (x-1)(x+1) = x-1 \Rightarrow x + 1 = 1 \Rightarrow x = 0.
  \]
  L'errore: si divide per $x - 1$, che vale $0$.
}

\ex{Il mentitore}{
  ``Io sono un bugiardo.'' Sono bugiardo sse quello che dico non è vero, sse
  ``io sono un bugiardo'' non è vero.
}

\ex{Paradosso di Berry}{
  ``Sia $x$ il più piccolo numero che non può essere definito con meno di 1000
  parole.'' Allora $x$ è definito (con meno di 1000 parole, quelle della frase)
  sse $x$ non è definibile con meno di 1000 parole.
}

\dfn{Metalinguaggio}{
  Un \textbf{metalinguaggio} è un linguaggio usato per parlare di un
  linguaggio. ``Metalinguistico'': relativo a un metalinguaggio.
}

\mprop{Perché nascono questi paradossi}{
  \begin{enumerate}
    \item uso \textbf{metalinguistico} del linguaggio naturale (la frase parla
          di frasi, di definizioni, \dots);
    \item \textbf{auto-applicazione} di un concetto metalinguistico (la frase
          parla di sé stessa).
  \end{enumerate}
  Rimedio: abbandonare il linguaggio naturale e usare un linguaggio
  opportuno della logica matematica, in cui linguaggio oggetto e
  metalinguaggio sono separati.
}

\subsubsection{Paradossi in matematica}

\thm{Paradosso di Russell}{
  Sia $X = \{ Y \mid Y \notin Y \}$. Allora $X \in X \iff X \notin X$.
}

\nt{
  \textbf{Soluzione:} $\{Y \mid Y \notin Y\}$ \emph{non è un insieme}.
  \begin{itemize}
    \item Non si può usare la nozione di insieme in modo metalinguistico: la
          collezione di tutti gli insiemi non è un insieme.
    \item Gli insiemi si costruiscono solo \textbf{selezionando elementi da un
          insieme già dato} (assioma di separazione di ZF:
          $\{ Y \in A \mid \varphi(Y)\}$). Così
          $\{Y \in \text{classe di tutti gli insiemi} \mid Y \notin Y\}$ è una
          \textbf{classe}, non un insieme, e quindi non contiene sé stessa.
  \end{itemize}
}

\subsubsection{Paradossi in informatica}

Molti linguaggi ammettono funzioni di ordine superiore (che prendono funzioni
in input): sono caratteristiche metalinguistiche che rendono possibile
l'auto-applicazione. Qui i paradossi \textbf{non sono evitabili}; ma invece di
produrre contraddizioni mostrano i \emph{limiti} dei nostri modelli di calcolo
(funzioni parziali, indecidibilità: vedi i Limiti 1--3).

\subsubsection{Il paradosso di Curry}

\ex{Versione in linguaggio naturale}{
  $A$ = ``Se questa frase è vera, allora Babbo Natale esiste''; sia $B$ =
  ``Babbo Natale esiste''.
  Per dimostrare $A$ (un'implicazione) assumiamo la premessa, ``questa frase
  è vera'', cioè $A$. Ma $A$ è $A \to B$, quindi per modus ponens $B$. Abbiamo
  dimostrato $A \to B$, cioè $A$. Ora $A$ è vera e $A \to B$ è vera: per modus
  ponens $B$. Babbo Natale esiste.
}

\mprop{Curry in un contesto logico}{
  \begin{enumerate}
    \item $X := (X \to Y)$ \hfill assunzione (``se questa frase è vera allora $Y$'')
    \item $X \to X$ \hfill legge di identità
    \item $X \to (X \to Y)$ \hfill da 2, sostituendo il lato destro con 1
    \item $X \to Y$ \hfill da 3 per \emph{contrazione}
    \item $X$ \hfill da 4, per la definizione 1
    \item $Y$ \hfill da 5 e 4 per modus ponens
  \end{enumerate}
}

\nt{
  \textbf{È di nuovo un paradosso?} Dipende dal passo 1. In molte logiche,
  se $Y$ è falso, la 1 non è possibile.
  In logica proposizionale classica (approfondimento): se $Y$ è falso,
  $X \leftrightarrow (X \to Y)$ diventa $X \leftrightarrow \neg X$, che non è
  soddisfatta da nessuna valutazione. Quindi non esiste alcuna formula $X$ con
  $X \equiv (X\to Y)$ e $Y$ falso: la ``definizione'' 1 non è esprimibile, e
  la derivazione è corretta ma parte da un'ipotesi impossibile.
  Il problema nasce solo se il linguaggio permette a una frase di riferirsi a
  sé stessa.

  (Sulla slide la derivazione è barrata a mano dal prof.)
}

% =====================================================================
\section{Perché studiare logica in AI}
% =====================================================================

\begin{enumerate}[label=\alph*)]
  \item È uno strumento per \textbf{esprimere il ragionamento umano}.
  \item È il fondamento di un \textbf{approccio classico all'AI} (vedi anche
        il corso Foundations).
  \item È il fondamento dei \textbf{linguaggi (dichiarativi) usati in AI}.
  \item Ha applicazione diretta in linguaggi usati in AI, come
        \textbf{Prolog e CLP} (questo corso).
\end{enumerate}

% =====================================================================
\section{Che cos'è l'AI: un po' di storia}
% =====================================================================

Come l'informatica, anche l'AI esiste \emph{prima} dei computer.

\begin{itemize}
  \item \textbf{Ramon Llull} (1232--1316), \emph{Ars combinatoria}: dischi di
        carta che combinano concetti rappresentati simbolicamente.
  \item \textbf{Cartesio} (1596--1650): influenzato dagli automi esposti a
        Parigi, studia il rapporto mente--corpo (\textbf{dualismo}: il corpo
        funziona come una macchina, la mente non è materiale, interagiscono
        nella ghiandola pineale). Precursore delle scienze cognitive e
        dell'AI perché considera la costruzione di automi un processo
        cognitivo importante.
  \item \textbf{Leibniz} (1646--1716), \emph{calculus ratiocinator}: un
        linguaggio universale in cui risolvere ogni controversia calcolando
        (\emph{``calculemus!''}).
  \item XVIII secolo: automi di \textbf{Vaucanson} (l'anatra che ``mangia'') e
        \textbf{Jaquet-Droz} (lo scrivano e il musicista, con movimenti
        controllati da un \emph{disco metallico}: un programma ante litteram).
  \item Il \textbf{Turco meccanico} (XVIII sec.): un ``automa'' che giocava a
        scacchi, in realtà con un giocatore umano nascosto. \emph{AI ante
        litteram\dots\ fake!} (Il nome è stato ripreso da Amazon Mechanical
        Turk, lavoro umano nascosto dietro servizi ``automatici''.)
  \item \textbf{Turing}: 1935 macchina universale; 1947 talk ``What we want is
        a machine that can learn from experience''; 1948 \emph{Intelligent
        Machinery} (mai pubblicato in vita, un vero manifesto dell'AI con
        sistemi simbolici e reti neurali); 1950 \emph{Computing Machinery and
        Intelligence} (imitation game).
  \item \textbf{Neuroni e apprendimento}: McCulloch--Pitts (1943), primo
        modello computazionale del neurone (logica a soglia); Hebb (1949),
        apprendimento hebbiano; Rosenblatt (1958), \textbf{perceptron}.
  \item \textbf{1956, Dartmouth}: John McCarthy conia il termine
        \emph{Artificial Intelligence}: ``ogni aspetto dell'apprendimento o
        dell'intelligenza può in principio essere descritto così precisamente
        da poter costruire una macchina che lo simuli''. McCarthy inventa
        anche il LISP.
\end{itemize}

\subsection{Le quattro visioni dell'AI (Russell--Norvig)}

\begin{center}
\begin{tabular}{l|l}
  \textbf{Thinking humanly} & \textbf{Thinking rationally} \\ \hline
  \textbf{Acting humanly} & \textbf{Acting rationally}
\end{tabular}
\end{center}

\dfn{Thinking humanly}{
  Sistemi che \textbf{pensano come gli umani}: imitare il cervello.
  Non ne sappiamo abbastanza (Broca 1861, Golgi 1873, EEG 1929, fMRI 1990):
  conosciamo relazioni tra aree e funzioni, non come funzionano un processo
  cognitivo o la memoria.
  L'approccio delle \textbf{scienze cognitive} (anni '60, ``rivoluzione
  cognitiva'' contro il comportamentismo) richiede teorie validate o
  top-down (predire e testare il comportamento di soggetti umani) o
  bottom-up (dati neurologici). Oggi è distinto dall'AI.
}

\dfn{Thinking rationally}{
  Sistemi che \textbf{pensano razionalmente}: le ``leggi del pensiero'',
  cioè la \textbf{logica}. Approccio \emph{normativo} (prescrittivo) e non
  descrittivo, da Aristotele in poi.
  Problemi: (1) non tutto il comportamento intelligente passa per la
  deliberazione logica; (2) qual è lo scopo del pensare, quali pensieri
  \emph{dovrei} avere?
}

\dfn{Acting rationally}{
  Sistemi che \textbf{agiscono razionalmente}: fare la cosa giusta, cioè
  quella che ci si aspetta massimizzi il raggiungimento dell'obiettivo date le
  informazioni disponibili. Non richiede necessariamente il pensiero (es.\
  il riflesso di chiudere gli occhi), ma il pensiero dovrebbe essere al
  servizio dell'azione razionale.
}

\dfn{Acting humanly: il test di Turing}{
  Turing (1950) sostituisce ``le macchine possono pensare?'' con ``le
  macchine possono \emph{comportarsi} in modo intelligente?'', e propone un
  test operativo: l'\textbf{imitation game}. Un interrogatore dialoga per
  scritto con un umano e una macchina e deve capire chi è chi.
}

\nt{
  Turing prevedeva che entro il 2000 una macchina avrebbe avuto il 30\% di
  probabilità di ingannare una persona comune per 5 minuti; anticipò le
  principali obiezioni all'AI e indicò le componenti chiave: conoscenza,
  ragionamento, linguaggio, apprendimento.
  \textbf{Problema}: il test non è riproducibile, costruttivo, né analizzabile
  matematicamente.

  Sul ``solipsismo'' Turing osserva che l'unico modo di sapere se una
  macchina pensa sarebbe \emph{essere} la macchina; ma lo stesso varrebbe per
  gli altri esseri umani. Invece di discutere all'infinito, ``è usuale avere
  la cortese convenzione che tutti pensino''.
}

\nt{
  \textbf{Posizione del corso:} da informatici adottiamo il punto di vista
  ``semplice'' di Turing, e consideriamo AI i sistemi che agiscono
  umanamente (o almeno razionalmente).
}

\dfn{Definizione ufficiale di AI}{
  L'AI è la capacità di un computer di svolgere compiti comunemente associati
  a esseri intelligenti. Due approcci principali:
  \begin{enumerate}
    \item \textbf{AI simbolica};
    \item \textbf{AI sub-simbolica} (connessionista, incluso il ML).
  \end{enumerate}
}

\nt{
  Peter Norvig: l'AI è \emph{``knowing what to do when you don't know what to
  do''}.
}

% =====================================================================
\section{AI simbolica vs.\ sub-simbolica}
% =====================================================================

\subsection{AI simbolica: i sistemi esperti}

\ex{Un sistema per le infezioni batteriche (stile MYCIN)}{
  \textbf{Sys}: la coltura è Gram positiva? \textbf{User}: NO.\\
  \textbf{Sys}: è Gram negativa? \textbf{User}: SÌ.\\
  \textbf{Sys}: il paziente è allergico agli antibiotici? \textbf{User}: NO.\\
  \textbf{Sys}: suggerirei una terapia a base di antibiotici.

  \smallskip
  Non è ChatGPT: il sistema \textbf{spiega} le proprie decisioni.
  Alla domanda \textbf{WHY?} (perché ti serve questa informazione?) risponde
  indicando la regola che sta cercando di applicare; alla domanda
  \textbf{HOW?} (come hai dedotto questo?) risponde indicando la regola usata
  e le premesse.
}

\dfn{Sistema esperto}{
  Un \textbf{sistema esperto} (anni '70 in poi):
  \begin{itemize}
    \item usa conoscenza di dominio fornita da umani, tipicamente come
          \textbf{regole logiche} (\textbf{knowledge base});
    \item risolve problemi in un dominio limitato con prestazioni simili a un
          esperto;
    \item cerca e costruisce dinamicamente la soluzione con un
          \textbf{motore inferenziale} che elabora le regole.
  \end{itemize}
}

\mprop{Caratteristiche della computazione simbolica}{
  \begin{enumerate}[label=\Roman*.]
    \item ragionamento basato sulla logica;
    \item può fornire \textbf{spiegazioni} delle decisioni;
    \item non adatta a problemi come il riconoscimento di immagini (che non si
          risolvono algoritmicamente).
  \end{enumerate}
}

\subsubsection{Formalizzare la diagnosi in logica}

\textbf{Obiettivo}: suggerire un farmaco adatto al paziente, cioè dimostrare
$\mathit{prescribe}(\mathit{Drug})$.

\textbf{Fatti:} $\mathit{gram}(\mathit{neg})$, \
$\neg\mathit{allergic}(\mathit{antb})$.

\textbf{Regole:}
\begin{itemize}
  \item[R1:] $\mathit{gram}(\mathit{neg}) \to \mathit{id}(\mathit{ecoli})$
             \hfill (Gram-negativo $\Rightarrow$ enterium-coli)
  \item[R2:] $\mathit{gram}(\mathit{pos}) \to \mathit{id}(\mathit{strep})$
             \hfill (Gram-positivo $\Rightarrow$ streptococco)
  \item[R3:] $\mathit{id}(\mathit{strep}) \vee \mathit{id}(\mathit{bact}) \to \mathit{ind}(\mathit{pen})$
             \hfill (strepto o batterio $\Rightarrow$ indicata la penicillina)
  \item[R4:] $\mathit{ind}(X) \wedge \neg\mathit{allergic}(X) \to \mathit{prescribe}(X)$
  \item[R5:] $\mathit{id}(\mathit{ecoli}) \to \mathit{ind}(\mathit{antb})$
\end{itemize}

\dfn{Motore inferenziale (backward chaining)}{
  Si parte dal \textbf{goal} e si va all'indietro:
  \begin{enumerate}
    \item il goal $\mathit{prescribe}(\mathit{Drug})$ fa \emph{match} con la
          testa di R4 ($\mathit{Drug}/X$); si attiva R4: nuovi sotto-goal
          $\mathit{ind}(X)$ e $\neg\mathit{allergic}(X)$;
    \item $\mathit{ind}(X)$ fa match con R3 ($X/\mathit{pen}$) e con R5
          ($X/\mathit{antb}$): due rami alternativi;
    \item ramo R3: serve $\mathit{id}(\mathit{strep})$ (via R2, serve
          $\mathit{gram}(\mathit{pos})$: \textbf{NO}) oppure
          $\mathit{id}(\mathit{bact})$ (nessuna regola: \textbf{NO}).
          Il ramo fallisce;
    \item ramo R5: serve $\mathit{id}(\mathit{ecoli})$, via R1 serve
          $\mathit{gram}(\mathit{neg})$: è un fatto, \textbf{SÌ};
          con $X = \mathit{antb}$ anche $\neg\mathit{allergic}(\mathit{antb})$
          è un fatto: \textbf{SÌ}.
  \end{enumerate}
  Risposta: $\mathit{prescribe}(\mathit{antb})$, con il legame
  $X/\mathit{antb}$ (annotato a mano dal prof sull'albero).
}

\begin{note}
  \textbf{Approfondimento: lo stesso in Prolog.} È esattamente il modo di
  ragionare di Prolog (risoluzione SLD: goal-directed, depth-first,
  sinistra-destra, backtracking):
\begin{lstlisting}[language=Prolog]
gram(neg).
not_allergic(antb).

id(ecoli)  :- gram(neg).          % R1
id(strep)  :- gram(pos).          % R2
ind(pen)   :- id(strep).          % R3 (due clausole per l'OR)
ind(pen)   :- id(bact).
ind(antb)  :- id(ecoli).          % R5
prescribe(X) :- ind(X), not_allergic(X).   % R4

?- prescribe(Drug).
Drug = antb.
\end{lstlisting}
  Osservazioni: la disgiunzione nel corpo di R3 diventa \emph{due clausole};
  il fatto negativo $\neg\mathit{allergic}$ è stato reso con un predicato
  positivo, perché le clausole di Prolog (clausole di Horn) non ammettono
  negazioni nei fatti. Le spiegazioni WHY/HOW corrispondono a mostrare il
  ramo dell'albero di derivazione.
\end{note}

\subsection{Logica e AI}
\begin{itemize}
  \item Molte logiche: proposizionale, primo ordine, modali (epistemica,
        temporale, deontica), \dots
  \item Molte tecniche di ragionamento: \textbf{deduzione}, \textbf{induzione},
        \textbf{abduzione}.
  \item Molti linguaggi e strumenti: Prolog, constraint programming.
  \item Base di molti sistemi esperti.
\end{itemize}

\nt{
  \textbf{Deduzione, induzione, abduzione} (approfondimento). Con regola
  $A \to B$:
  \begin{itemize}
    \item \emph{deduzione}: da $A$ e $A \to B$ concludo $B$ (corretta);
    \item \emph{abduzione}: da $B$ e $A \to B$ ipotizzo $A$ come spiegazione
          (es.\ diagnosi: dal sintomo alla causa; non corretta in generale);
    \item \emph{induzione}: da molti casi di $A$ e $B$ insieme generalizzo
          $A \to B$ (è ciò che fa il ML).
  \end{itemize}
}

\subsection{AI sub-simbolica}

\mprop{Caratteristiche della computazione sub-simbolica}{
  \begin{enumerate}[label=\Roman*.]
    \item basata su reti neurali (artificiali) e ML;
    \item molto buona per riconoscimento di immagini e problemi non
          risolvibili algoritmicamente;
    \item \textbf{non può fornire spiegazioni} delle decisioni (e quindi è
          ancora lontana dal superare il test di Turing).
  \end{enumerate}
}

\ex{Riconoscimento di tumori in istologia}{
  \textbf{Train}: da immagini di vetrini interi si estraggono patch etichettate
  (tumore / normale) e si addestra un modello profondo.
  \textbf{Test}: il vetrino nuovo viene diviso in patch sovrapposte, il modello
  restituisce una probabilità per patch, e si ottiene una \emph{mappa di
  probabilità} del tumore. Il modello indica \emph{dove}, ma non \emph{perché}.
}
